\section{Exact remaining gates}
\status{Open; this document is an incomplete research return.}
For A.1 and L1, our completed gluing and local-to-global propositions require
actual representing schemes over the local bases. A uniform Artin--Rees
argument has now passed a targeted proof check: conditional on those local
representations, condition~8 on finite-type tests yields openness of the
\emph{\'etaleness} locus, hence an algebraic space. The separated \'etale
relation has finite fibres over a Noetherian scheme and is quasi-compact;
valuative uniqueness then gives a separated, locally quasi-finite scheme.
The remaining main gate is deriving the needed local schemes from the literal
actual-modularity premise of condition~7, rather than a stronger
pro-modularity premise. We do not identify these premises.

L2 remains undecided. The dimension-one independence example settles the
other question of A.2. A henselized infinitesimal-group candidate has been
checked on complete local tests, but full fpqc descent or preservation of
filtered colimits after sheafification remains unproved; it is not a
counterexample in this paper.

B.3 has explicit failures on reduced, non-reduced, and non-Noetherian bases
and exact positive classes from the opened sources. We have not established
that their union is a largest possible class in the literature.

C.1 is reduced beyond regular Noetherian rings to rings whose modular fibres
can be singular. The arbitrary-ring obstruction sequence and explicit
normal-domain and torsion counterexamples are stated in Part~I. A necessary
and sufficient intrinsic classification for every singular mixed-characteristic
ring remains open in this run. For C.3, the representation-theoretic formula
over fields is proved; explicit unrestricted modular indecomposable
classification and the corresponding non-field computation remain.
For C.4, the noncommutative unipotent formal-group alternative and realization
without a baseline group remain; the proposed universal splitting statements
and the identification with Hall isoclinism are already refuted.

The component notes record attempted approaches and exact failure points.
The dated progress ledger distinguishes working updates from the first
formal review, due only after three hours of work.
